\section{Source audit and precise corrections}
\label{powers:sec:audit}

\subsection{Scope of the repository audit}

The source baseline is ProveIt commit
\texttt{b7667ae40a2957eed9a53d8559d53a3ee20dd1c1}, including the
current manuscript and the incoming archives inspected for this
continuation \cite{Repo}. The review concentrated on the polynomially
weighted resolvent, the multiple-Gamma normalization, continuous
resolvent powers, collision finite parts, and the quartic
power-to-harmonic bridge \cite{Quartic,Barnes}. No new false theorem
was identified in the selected repository statements used here.
This is a targeted mathematical audit, not a certification of every
chapter of the complete manuscript. The two corrections below concern
precisely identified external sources.

The finite-part convention remains the original unit coordinate:
\(\FP\int_0^1 f(x)\,\dd x\) retains the constant coefficient of
\(\int_\epsilon^1 f(x)\,\dd x\) after its explicit local power and
logarithmic terms are separated. It is not an unspecified analytic
regularization. Likewise, a constant coefficient in a spectral Laurent
series or in a polynomial asymptotic expansion in \(\log N\) is a
different operation until a comparison theorem proves equality. In the
present argument the normalized contour formula and finite summation
by parts provide that comparison. Matching principal parts alone would
leave an undetermined entire remainder and would not prove the moment
identity.

\subsection{A reversed pair of integrals in a Tornheim source}

The author-hosted version of Bailey and Borwein's
\emph{Computation and theory of Mordell--Tornheim--Witten sums II},
dated 5 May 2014, contains a sign reversal in equation~(110), on printed
page~24 \cite{BB2014}. Its preceding equation~(109) is correct. The
correction is recorded specifically for this inspected dated version;
no claim is made here about an independently uninspected final
typesetting.

For \(\Re a>0\), put
\begin{align}
 A(a)&=\frac1{\Gamma(a+1)}\int_0^1
       \frac{(-\log u)^a[-\log(1-u)]}{1-u}\,\dd u,
       \label{powers:audit:A}\\
 B(a)&=\int_0^1\Li_a(u)
                  \frac{-\log(1-u)}u\,\dd u.
       \label{powers:audit:B}
\end{align}
The power of \(-\log u>0\) uses its real logarithm. Both integrals
converge absolutely and depend holomorphically on \(a\) in this
half-plane.

\begin{proposition}[Corrected integral identity]
\label{powers:audit:BB-correction}
For \(\Re a>0\),
\begin{equation}
 \boxed{\displaystyle
 \zeta(a+2)=B(a)-A(a).}
 \label{powers:audit:correct-order}
\end{equation}
Equation~(110) of the specified source prints \(A(a)-B(a)\).
\end{proposition}

\begin{proof}
Write \(H_m=\sum_{j=1}^m j^{-1}\) and
\(T(a,1,1)=\sum_{m,n\ge1}m^{-a}n^{-1}(m+n)^{-1}\).
For each positive integer \(m\), the exact finite calculation
\[
 \sum_{n=1}^N\frac1{n(m+n)}
 =\frac{H_N+H_m-H_{N+m}}m
\]
gives \(\sum_{n\ge1}1/[n(m+n)]=H_m/m\). Therefore
\begin{equation}
 T(a,1,1)=\sum_{m\ge1}\frac{H_m}{m^{a+1}}
         =\zeta(a+1,1)+\zeta(a+2),
 \label{powers:audit:direct-T}
\end{equation}
where \(\zeta(s,1)=\sum_{m\ge1}H_{m-1}m^{-s}\).
These series converge absolutely when \(\Re a>0\).
Expanding the polylogarithm and logarithm in
\eqref{powers:audit:B} gives \(B(a)=T(a,1,1)\), while
\[
 \frac{-\log(1-u)}{1-u}=\sum_{n\ge1}H_nu^n,
 \qquad
 \int_0^1u^n(-\log u)^a\,\dd u
       =\frac{\Gamma(a+1)}{(n+1)^{a+1}}
\]
give \(A(a)=\zeta(a+1,1)\). Absolute convergence justifies all
interchanges, proving \eqref{powers:audit:correct-order}.
\end{proof}

There is an exact elementary counterexample to the printed sign.
At \(a=1\), substitution followed by the geometric series gives
\[
 A(1)=\int_0^1\frac{\log v\log(1-v)}v\,\dd v=\zeta(3),
 \qquad
 B(1)=\int_0^1\frac{\log^2v}{1-v}\,\dd v=2\zeta(3).
\]
Thus its printed right side is \(-\zeta(3)\), whereas its left
side is \(\zeta(3)>0\). This corrects the integral line; it does
not invalidate equation~(109) or the spectral derivative identities
obtained from that correct equation. The proof above also establishes
the whole half-plane identity directly, without a Carlson-theorem
argument.

\subsection{A distinction between multiple and ordinary zeta reductions}

On printed page~205 of \emph{Nonlinear Euler sums}, Mez\H{o}
states that it is not known whether all the nonlinear sums under
discussion can be expressed as multiple zeta values \cite{Mezo2014}.
For convergent products of positive-integer generalized harmonic
numbers, this target class is too broad: a finite multiple-zeta
reduction follows directly from ordering their summation indices.

To make the correction precise, use the strict finite sums
\[
 H_n(b_1,\ldots,b_\ell)
 =\sum_{n\ge j_1>\cdots>j_\ell\ge1}
       j_1^{-b_1}\cdots j_\ell^{-b_\ell}.
\]
For positive integers \(a_1,\ldots,a_d\), expand
\(\prod_{i=1}^d H_n^{(a_i)}\) by the ordered blocks of equal
indices. Every ordered set partition
\((B_1,\ldots,B_\ell)\) of \(\{1,\ldots,d\}\) contributes
\(H_n(b_1,\ldots,b_\ell)\), with
\(b_j=\sum_{i\in B_j}a_i\). This is the finite stuffle product.
For \(q\ge2\), each resulting term satisfies
\begin{equation}
 \sum_{n\ge1}\frac{H_n(b_1,\ldots,b_\ell)}{n^q}
 =\zeta(q,b_1,\ldots,b_\ell)
       +\zeta(q+b_1,b_2,\ldots,b_\ell).
 \label{powers:audit:finite-stuffle}
\end{equation}
The two terms distinguish \(n>j_1\) and \(n=j_1\).
All sums are absolutely convergent; hence this is a finite exact
reduction for every such nonlinear Euler sum.

For example,
\begin{align}
 \sum_{n\ge1}\frac{H_n^2}{n^2}
 &=2\zeta(2,1,1)+2\zeta(3,1)+\zeta(2,2)+\zeta(4)
   =\frac{17}{4}\zeta(4).
 \label{powers:audit:example}
\end{align}
The first equality already demonstrates the required finite
multiple-zeta reduction; the last equality uses the familiar
weight-four evaluations. A precise replacement for the source's
wording should identify a smaller target, such as ordinary zeta
values and their products, or a specified class of linear Euler
sums. Formula~\eqref{powers:audit:finite-stuffle} does not establish
either stronger reduction in general. It is a classical algebraic
mechanism, not a new arithmetic-independence result.

\subsection{Unchanged conjecture status}

The fifth-power identity and the algorithm for arbitrary integer
powers establish the finite harmonic spectral formulas stated in
this article. They do not prove that every retained spectral
derivative reduces to ordinary constants. A nonzero coordinate in a
specified formal presentation also does not establish arithmetic
independence of its numerical value. The current Gaussian
\(S_6\) and revised \(S_8\) identities in the pinned manuscript
remain conjectures \cite{Repo}; none of the arguments here changes
their status.
